% =====================================================================
% Chapter 2: Related Work
%
% Required topic: a critical review of the methods that exist at the
% time of writing. Every section describes what a line of work does
% and then where it falls short for global robustness, and the last
% section collects these limitations into the gap this thesis closes.
%
% Adapted from the Related Work section of the paper, whose paragraphs
% are Global Robustness Verification, Incremental Verification and
% Relaxations, and extended with two sections a thesis has room for,
% on attack-guided verification and on the propagation of differences.
%
% The entries are in references.bib and aaai2027.bib, copied from the
% paper.
% =====================================================================
\chapter{Related Work}
\label{ch:related-work}

In this chapter, we review the methods that verify the robustness of
neural networks. We begin with the verifiers of global robustness,
continue with the relaxations that scaled local robustness and
with the reuse of one verification for another, and close with the
propagation of intervals and of differences between networks and with
the attacks that guide a verifier. 

\section{Global Robustness Verification}
\label{sec:related-global-robustness}

Local robustness verification certifies that a single given input
cannot be misclassified under any small
perturbation~\citep{katz2017reluplex,gehr2018,singh2019,wang2021}. A
deployed network, however, faces inputs it has never seen, and a
guarantee that covers a fixed set of inputs says nothing about them.
Global robustness covers all of them: it guarantees that any input
classified above a confidence level cannot be flipped by a
perturbation~\citep{leino2021gloro,kabaha2024vhagar,athavale2024}, so
that at inference time any input below that level is flagged as
potentially vulnerable. The guarantee is stronger, and so is the cost
of obtaining it, which is why the verifiers below settle for a weaker
guarantee, for a restricted class of networks, or for a long analysis.

Several verifiers settle for a probabilistic guarantee.
\citet{blohm2025} certify a probabilistic relaxation of global
robustness by sampling an $\epsilon$-net and invoking a local
robustness oracle on each sample, and gRoMA measures probabilistic
global robustness by aggregating local statistical
scores~\citep{levy2023}. Both report how likely an input is to be
attackable rather than whether it is, so neither rules out an attack
on an input that its samples did not cover, and the cost of both grows
with the number of samples they draw.

Among the deterministic verifiers, DeepGlobal enumerates decision
boundaries by symbolic verification to identify adversarial regions,
but it requires a custom activation function, so it does not apply to
a network trained in the usual way~\citep{sun2022deepglobal}.
\citet{athavale2024} verify global robustness by self-composition and
a softmax abstraction, which yields sound but not tight bounds, and a
bound that is not tight flags inputs as vulnerable that an exact
analysis would clear.
% TODO (Dana): The AI review of the AAAI submission asked for sharper positioning against \citet{athavale2024}. Add one sentence on how their property and guarantee differ from ours (e.g., checking a given confidence level vs.\ computing the minimal bound, softmax abstraction vs.\ exact encoding), after checking their paper.
 \baseline{} computes the minimal globally robust
bound exactly, by mixed-integer programming, dependency analysis and a
hyper-attack, and it is an anytime verifier, that is, it returns the
bounds it has proved so far when it is stopped~\citep{kabaha2024vhagar}.
Exactness is what makes it expensive: its analysis takes hours on
networks of up to 2000 neurons, which leaves the larger networks out
of reach.

Other works address formulations close to global robustness.
Networks can be trained to be globally robust by
construction~\citep{leino2021gloro}, which changes the network rather
than analyzes a given one. \citet{wang2022} bound the maximal variation of the output under a perturbation over all inputs,
and \citet{tobler2025} verify the implementation of a
certifier rather than a network. The property this thesis verifies is
the one of \baseline.

\section{Relaxations}
\label{sec:related-relaxations}

An exact encoding of the problem
% TODO (Dana): Neta: which problem? Name it, e.g., local or global robustness verification.
 costs a binary variable per
undecided neuron
% TODO (Dana): Neta: ``undecided'' is a new term. Use ``unstable'', as in Chapter~\ref{ch:background}.
, so the search space grows exponentially with the number of unstable neurons. To
overcome this, many verifiers use relaxations, which replace the
binary variable that encodes a neuron exactly with a convex region
that constrains it. The triangle relaxation
is the standard one, bounding a ReLU by the smallest triangle over its
pre-activation interval~\citep{ehlers2017,singh2019}, and the zonotope
abstraction propagates linear expressions over the perturbation, which
keeps the dependencies between neurons that intervals
lose~\citep{singh2018fast,gehr2018,mirman2018differentiable}. Later
work tightens the relaxation by considering several neurons at once,
as in $k$-ReLU~\citep{singh2019krelu} and
PRIMA~\citep{muller2022prima}, by partitioning the input
region~\citep{anderson2020tightened}, by relaxing a ReLU together with
its pre-activation~\citep{tjandraatmadja2020convex}, and by dual
algorithms over sets of active
variables~\citep{depalma2021}. \citet{wang2022} over-approximate their encoding by decomposing the network and relaxing parts of it to a linear program, which they then refine.

However, except for \citet{wang2022}, none of these relaxations relies on the two copies $N$ and
$N^{p}$ of the global robustness setting, and on the relation between
$z_{m,k}$ and $z^{p}_{m,k}$, to obtain a tighter region.
\citet{wang2022} encode the two copies with extra dependencies between them, and over-approximate this encoding by network decomposition and a linear-program relaxation, so their bound is sound but not exact. \tool relaxes only the neurons whose phase is nearly determined, conditions the relaxation of a neuron in one copy on the phase of its counterpart in the other copy, and encodes the remaining neurons exactly.
%
% TODO (Dana): Neta, double-check this description of \citet{wang2022} (interleaving twin-network MILP, network decomposition, LP relaxation with refinement) against their paper. The AI review of the AAAI submission asked for this positioning.


\section{Incremental Verification}
\label{sec:related-reuse}

Neural networks in real-world systems are rarely static.
They undergo continuous adaptation through
retraining~\citep{parisi2019}, fine-tuning~\citep{yosinski2014},
pruning and quantization~\citep{han2016deepcompression}. Because
complete formal verification is
NP-complete~\citep{katz2017reluplex},
re-verifying an entire network from scratch after minor modifications
quickly becomes intractable.
Incremental verifiers answer it by carrying the analysis of
one network over to the next exploiting structural and behavioral invariants
across model revisions. They transfer proof
templates~\citep{ugare2022fanc,fischer2022}, reuse proofs between
solver calls~\citep{yang2023deepinc}, keep a continuous
train--verify--deploy loop~\citep{ugare2023continuous}, and reuse
probabilistic certificates~\citep{ugare2024irs}.

What they share is the property they carry over. All of them target
local robustness, or global properties that do not capture a
perturbation. IVAN verifies global properties
incrementally~\citep{ugare2023}, but in the scenario of robustness it
loses precision over perturbations. Therefore, there is still room to
research new approaches for incremental verification of global robustness
that do not lose precision.

\section{Interval and Difference Propagation}
\label{sec:related-difference-propagation}

Interval Bound Propagation provides a computationally efficient
abstraction for neural network verification,
serving as a foundational primitive in scalable
local robustness verifiers~\citep{liu2021algorithms}.
ReluVal carries linear expressions over the input instead of constants, which
keeps the dependencies that plain intervals lose~\citep{wang2018reluval},
but it too is defined over a single network and a single input.


Intervals serve global properties as well. Fairify verifies
individual fairness, a property of every pair of similar inputs, and
uses interval arithmetic to decide the neurons that stay inactive
over a part of the input domain~\citep{biswas2023fairify}, and global
specifications in general can be reduced to ordinary ones by
composing two copies of the network and handing the composition to an
existing verifier~\citep{boetius2023hyperproperties}.
%
% TODO (Dana): Neta: avoid orphan paragraphs. The first two paragraphs of this section are too short. Extend or merge them.


Differential verification propagates divergence bounds between
network pairs rather than analyzing each copy alone.
ReluDiff and NeuroDiff bound how far apart two different networks are on the
same input~\citep{paulsen2020reludiff,paulsen2020neurodiff}.
LUCID bounds the difference between networks to verify a privacy
property~\citep{kabaha2025lucid}. In all of them the two networks are
distinct and the input is shared. However, in global robustness settings,
we have two copies of the same network, with different inputs,
where the inputs differ by a perturbation.

Another common abstraction is Zonotopes. In local robustness they keep the
dependencies between neurons that intervals
lose~\citep{singh2018fast,gehr2018}. RefineZono refines the bounds of
selected neurons with a solver, so that the zonotope certifies
more~\citep{singh2019refinezono}. For global robustness,
\citet{athavale2026differential} propagate differential halo zonotopes
to decide a confidence-based two-safety formulation of
global robustness while losing precision.

\section{Attack-Guided Verification}
\label{sec:related-attack-guided}

An adversarial attack searches for a perturbation that changes the
classification of an input~\citep{szegedy2014,carlini2017,madry2018}.
An attack is incomplete: when it finds a perturbation, the input is
proved to be vulnerable, and when it fails, nothing is proved. Its
value to a verifier is therefore not the proof but the solution it
produces, which a constraint solver can start the search from, and
which bounds the answer from one side while the solver works on the
other.

Verifiers of global robustness use attacks in exactly this way.
\baseline searches for an attack over the two copies of the network,
a hyper-attack, and submits it to the solver together with the
encoding~\citep{kabaha2024vhagar}. The gain is real but it is bounded
by what an attack can be: it is one feasible solution, and it leaves
the encoding that the solver has to reason about untouched. The
abstractions of this thesis act on that encoding instead, and the two
are complementary, which is why \tool retains the hyper-attack of
\baseline.

% \section{Limitations of Existing Methods}
% \label{sec:related-gap}

% The review above leaves four limitations, and together they describe
% the state of global robustness verification at the time of writing.
% First, the verifiers that scale do so by weakening the guarantee, to a
% probabilistic one or to a sound but loose bound, and a loose bound
% flags inputs as vulnerable that an exact analysis would clear. Second,
% the one verifier that is exact takes hours on networks of up to 2000
% neurons. Third, the abstractions that made local robustness tractable
% are defined over a single network, and the works that propagate
% differences do so between two distinct networks on a common input,
% rather than between a network and its own perturbed copy. Fourth, the
% reuse of a verification has been shown for local robustness and for
% perturbation-free global properties, so a revised network has to be
% verified from scratch.

% The purpose of this thesis follows from these four limitations.
